\documentclass[11pt]{article}
\usepackage[utf8]{inputenc}
\usepackage[T1]{fontenc}
\usepackage[margin=2cm]{geometry}
\usepackage{tikz}
\usepackage{tcolorbox}
\usepackage{array}
\usepackage{amsmath}
\usepackage{amssymb}
\usepackage{graphicx}
\usepackage{fancyhdr}
\usepackage{enumitem}

\tcbuselibrary{skins}
%\pagestyle{empty}

% ---- En-tête Nom / Groupe ----
\newcommand{\entete}{%
  \noindent Nom: \rule{6.5cm}{0.4pt} \hfill Groupe: \rule{3cm}{0.4pt}\\[4mm]
}

% ---- Boîte "Cible" avec titre en gras dans la boîte ----
\newtcolorbox{cible}[1]{
  colback=white, colframe=black, boxrule=0.6pt, arc=0pt,
  left=4mm, right=4mm, top=2mm, bottom=2mm,
  fonttitle=\bfseries, coltitle=black,
  before upper={\textbf{#1}\par\smallskip}
}

% ---- Case blanche générique (pour les chaînes d'opérations) ----
%\newcommand{\caseop}[1]{%
%  \begin{minipage}[t][3cm][t]{\linewidth}
%  #1
%  \end{minipage}
%}

\begin{document}

%==================== PAGE 1 ====================
\entete
\noindent \textbf{Évaluation: Les fractions} \hfill Chapitre 1
\vspace{2mm}
\hrule
\vspace{4mm}

\begin{cible}{Cible 1: Lire, écrire et représenter une fraction et un nombre fractionnaire. Rechercher une fraction équivalente.}

1. Donne la fraction représentée dans chaque cas.  

\vspace{2mm}

\begin{tabular}{@{}p{0.5cm}p{2.3cm}p{4.3cm}p{0.5cm}p{2.3cm}p{4.3cm}@{}}
a) & $2^3 = $ & \rule{3.2cm}{0.4pt} & b) & $4^2 = $ & \rule{3.2cm}{0.4pt} \\[4mm]
c) & $9^0 = $ & \rule{3.2cm}{0.4pt} & d) & $8^1 = $ & \rule{3.2cm}{0.4pt} \\[4mm]
\end{tabular}

\vspace{6mm}

\noindent 2. Rempli le tableau suivant : 


\vspace{6mm}
\noindent 3. Transforme les fractions suivantes en nombre fractionnaire.

\vspace{2.5cm}


\vspace{6mm}

\noindent 4. Transforme les nombre fractionnaires en fractions.

\vspace{2.5cm}

\vspace{4mm}
\end{cible}


\newpage

%==================== PAGE 2 ====================

\begin{cible}{Cible 1 (suite): Lire, écrire et représenter une fraction et un nombre fractionnaire. Rechercher une fraction équivalente.}

\noindent 5. Trouve la valeur manquante pour que les fractions soient équivalentes. Montre tes calculs. 

\begin{cible}{Cible 2: Situer, comparer et ordonner des fractions, des pourcentages et des entiers.}

\noindent 1. Place le symbole <, > ou =


\noindent 2. Place les nombres  suivants en ordre croissant (du plus petit au plus grand): 

\vspace{2mm}



\newpage

%==================== PAGE 3 ====================


\begin{cible}{Cible 3: Effectuer des chaînes d'opérations.}

1. Trouve le résultat de la chaine d'opérations suivante. Écris tes étapes de calcul.

\vspace{1cm}

\renewcommand{\arraystretch}{1}
\begin{tabular}{|p{7.7cm}|p{7.7cm}|}
\hline
\rule{0pt}{4mm} a) $6 \times (4+3) = $ & b) $7 - (-6)^2 = $ \\[8cm]
\hline
\rule{0pt}{4mm} c) $ (-2)^3 - 7 = $ & d) $5 \times -9 + 4 = $ \\[8cm]
\hline
\end{tabular}

\end{cible}

\newpage

\begin{cible}{Cible 3: Effectuer des chaînes d'opérations.(suite)}

1. Trouve le résultat de la chaine d'opérations suivante. Écris tes étapes de calcul.

\vspace{1cm}

\renewcommand{\arraystretch}{1}
\begin{tabular}{|p{7.7cm}|p{7.7cm}|}
\hline
\rule{0pt}{4mm} e) $5 \times 9 - 50 \div (12-2) = $ & f) $ 4 \times 6 - (-4)^3 = $ \\[8cm]
\hline
\rule{0pt}{4mm} g) $(-7+5)^3 \div 4 = $ & h) $(-7)^2 - 8 + 4 \div -4 = $ \\[8cm]
\hline
\end{tabular}

\end{cible}


\end{document}